

% !TEX root = ./ipeesec.tex
%%%%%%%%%%%%%%%%%%%%%%%%%%%%%%%%%%%%%%%%%%%%%%%%%%%%%%%%%%%
% Obere Titelmakros. Editieren Sie diese Datei nur, wenn
% Sie sich ABSOLUT sicher sind, was Sie da tun!!!
% Version 2.2 - 2.3.2010
%%%%%%%%%%%%%%%%%%%%%%%%%%%%%%%%%%%%%%%%%%%%%%%%%%%%%%%%%%%
\newcommand{\AN}{twoside}
\newcommand{\AUS}{}


%\newcommand{\englisch}{}
%\newcommand{\deutsch}{\usepackage[german]{babel}}

%% Die folgenden auskommentierten Optionen dienen der automatischen
%% Erkennung des Latex-Kompilers und dem Setzen der davon abhängigen
%% Einstellungen. Bei Problem z.B. mit dem Einbinden von verschiedenen
%% Grafiktypen bei Verwendung von PdfLatex oder Latex, einfach die
%% verschiedenen \usepackage(s) ausprobieren. (Mit diesen Einstellungen
%% funktionierte diese Vorlage bei der Verwenundg von latex.exe als
%% Kompiler bei den meisten Studierenden.)

%\newif\ifpdf \ifx\pdfoutput\undefined
%\pdffalse % we are not running pdflatex
%\else
%\pdfoutput=1 % we are running pdflatex
%\pdfcompresslevel=9 % compression level for text and image;
%\pdftrue \fi

\documentclass[11pt,a4paper, \zweiseitig]{article}
\usepackage{ifthen}
\usepackage{xstring} % \StrBehind, used for the bibliography URLs


%\usepackage[iso]{umlaute}
\usepackage{iftex}
\ifPDFTeX
	\PackageError{imports}{This document must be compiled with XeLaTeX (or LuaLaTeX)}{Run xelatex (see the recipe in .vscode/settings.json or .latexmkrc), not pdflatex.}
\fi
% inputenc/fontenc are for pdfLaTeX only; XeLaTeX reads UTF-8 natively and
% fontspec replaces the OT1/T1 font-encoding machinery.
\usepackage{fontspec}
\defaultfontfeatures{Ligatures=TeX}
\setmainfont{TeX Gyre Pagella} % OpenType clone of Palatino, replaces \usepackage{palatino}
% Monospace font used for \texttt and code listings (minted below). JuliaMono
% is designed for scientific/programming use and has by far the broadest
% Unicode coverage of any monospace font (math operators, arrows, blackboard
% bold, etc.), which is exactly the symbol set used in the pseudocode in this
% document (▶, ×, ℕ, Ω, λ, ∀, ≤, ...). This lets code blocks use those symbols
% as plain UTF-8 characters instead of dropping into \mathescape.
\setmonofont{JuliaMono}[Scale=MatchLowercase]
% Fallback font for symbols/emoji outside the main font (e.g. \emoji{🌳}).
% Symbola covers a very large swath of Unicode, incl. most emoji, in outline
% (monochrome) form, which is what XeLaTeX can actually typeset; the color
% "Noto Color Emoji" font will not render in color under XeLaTeX/LuaLaTeX.
\newfontfamily\emojifont{Symbola}
\newcommand{\emoji}[1]{{\emojifont #1}}
%\usepackage{makeidx} % um ein Index zu erstellen
\usepackage[nottoc]{tocbibind}
\usepackage{fancybox} % fuer die Rahmen
\usepackage{shortvrb}
\usepackage{url}
\usepackage{xcolor}
\usepackage{nicefrac}
% minted (https://ctan.org/pkg/minted) replaces `listings`: it renders code
% via Pygments, so any Unicode symbol supported by the chosen font (see
% \setmonofont above) can be typed directly in the listing, and it produces
% much nicer highlighted/framed output ("beautiful pseudocode") than
% `listings`. It shells out to Pygments (via the `latexminted` helper), so
% compilation needs -shell-escape (already set in .latexmkrc / .vscode).
\usepackage{minted}
\usemintedstyle{bw} % Pygments "bw": bold/italic only, no color or background
\setminted{
	fontsize=\small,
	breaklines=true,
	autogobble=true, % strip common leading whitespace, so indented \begin{minted} blocks aren't shifted right
	linenos=true, % classic look: line numbers in the margin,
	frame=leftline, % separated from the code by a vertical rule
	framesep=6pt,
	numbersep=8pt,
	ignorelexererrors=true, % Pygments flags non-ASCII like ▶ as errors and draws a red box around them
	baselinestretch=1.2,
}
\renewcommand{\theFancyVerbLine}{\scriptsize\arabic{FancyVerbLine}}
\setlength{\marginparwidth}{2cm}
\usepackage{todonotes}
\usepackage{enumitem} % (a), (b), ... labels in definitions
\usepackage[colorlinks,citecolor=black,linkcolor=black,urlcolor=black]{hyperref} %anklickbares Inhaltsverzeichnis

\ifthenelse{\equal{\biber}{true}}{
	% only needed for biber
	\usepackage{csquotes}
	\usepackage[style=numeric,sorting=none,natbib=true,backend=biber,mincitenames=1,maxcitenames=2,maxbibnames=99,uniquelist=false]{biblatex}

	% https://tex.stackexchange.com/a/334703/8850
	\AtEveryBibitem{%
		\clearfield{issn}
		\clearfield{isbn}
		\clearfield{doi}
		\clearfield{location}
		\clearlist{location}
		\clearlist{address}

		\ifentrytype{online}{}{% Remove url except for @online and @thesis
			\ifentrytype{thesis}{}{\clearfield{url}}
		}
	}
	% no colons in the bibliography: "In Proceedings" instead of "In: Proceedings",
	% and "URL openvpn.net/" (still linking to the full URL) instead of "URL: https://..."
	\renewbibmacro*{in:}{\bibstring{in}\space}
	\DeclareFieldFormat{url}{%
		\mkbibacro{URL}\addspace
		\StrBehind{#1}{://}[\urlwithoutscheme]%
		\href{#1}{\expandafter\nolinkurl\expandafter{\urlwithoutscheme}}}
}
{}%no else

% Falls es bei \citet ein Komma zwischen Name und Jahr gibt:
% https://tex.stackexchange.com/questions/312539/unwanted-comma-between-author-and-year-using-citet-command
% (thx @ Markus Brenneis)
%\DeclareDelimFormat[cbx@textcite]{nameyeardelim}{\addspace}



\ifthenelse{\equal{\sprache}{deutsch}}{
	\usepackage[ngerman]{babel}
	% Bibtex u.a -> et al.
	\ifthenelse{\equal{\biber}{true}}{
		\DefineBibliographyStrings{ngerman}{
			andothers = {{et\,al\adddot}},
		}
		\newcommand{\references}{Literatur}
	}
	{} % do nothing when not using biber
	\usepackage[autostyle, german=quotes]{csquotes} % Deutsche Anführungszeichen im Literaturverzeichnis (thx @ Markus Brenneis)

}{ \newcommand{\references}{References}}

\usepackage{a4wide} % ganze A4 Weite verwenden



%\ifpdf
%\usepackage[pdftex,xdvi]{graphicx}
%\usepackage{thumbpdf} %thumbs fuer Pdf
%\usepackage[pdfstartview=FitV]{hyperref} %anklickbares Inhaltsverzeichnis
%\else
%\usepackage[dvips,xdvi]{graphicx}
\usepackage{graphicx}
\usepackage{caption}
\captionsetup{labelsep=period} % "Figure 1." instead of "Figure 1:"
\usepackage{subcaption}
\usepackage[section]{placeins} % floats stay within their section
\usepackage{array}
\usepackage{booktabs}

%\fi

% Mathematik / kategorientheoretische Notation und kommutative Diagramme
\usepackage{amsmath}
\usepackage{mathtools} % \coloneqq
\usepackage{amssymb}
\usepackage{amsthm}
\usepackage{stmaryrd} % \llbracket, \rrbracket
\usepackage{tikz-cd}
\usepackage{tikz-3dplot}
\usepackage{pgfplots}
\usepackage[super]{nth}
\pgfplotsset{compat=1.18}
\usetikzlibrary{3d, positioning, arrows.meta, patterns, patterns.meta, decorations.pathreplacing}
% draft watermark: drawn diagonally on top of every page, in translucent red
\usepackage{eso-pic}
\AddToShipoutPictureFG{%
	\put(\LenToUnit{0.5\paperwidth},\LenToUnit{0.5\paperheight}){%
		\tikz[overlay]\node[rotate={atan(\paperheight/\paperwidth)}, text=red,
			opacity=0.2, font=\bfseries] at (0,0)
		{\resizebox{29cm}{!}{DRAFT (lots of errors!)}};}} % A4 diagonal is about 36cm
\usepackage{mathpartir} % inference rules for the internal logic
\usepackage{aliascnt}
\usepackage[capitalise,nameinlink]{cleveref}

\newtheorem{theorem}{Theorem}[section]
% each environment gets an alias of the theorem counter, so that the shared
% numbering is kept but cleveref sees the environment name, not "theorem"
\newcommand{\newsharedtheorem}[2]{%
	\newaliascnt{#1}{theorem}%
	\newtheorem{#1}[#1]{#2}%
	\aliascntresetthe{#1}}
\newsharedtheorem{definition}{Definition}
\newsharedtheorem{proposition}{Proposition}
\newsharedtheorem{lemma}{Lemma}
\newsharedtheorem{corollary}{Corollary}
\newsharedtheorem{example}{Example}
\newsharedtheorem{remark}{Remark}

\newcommand{\Set}{\mathbf{Set}}
\newcommand{\restr}[2]{#1|_{#2}}
\newcommand{\NN}{\mathbb{N}}
\newcommand{\ZZ}{\mathbb{Z}}
\newcommand{\List}{\mathrm{List}}
\newcommand{\Hom}{\text{Hom}}
\newcommand{\im}{\text{im }}
% semantic brackets for the interpretation of a specification
\newcommand{\sem}[1]{[\![#1]\!]}
\newcommand{\id}{\text{id}}
\newcommand{\true}{\text{true}}
\newcommand{\false}{\text{false}}
\newcommand{\Bool}{\text{Bool}}
% topos of trees and guarded recursion
\newcommand{\toposoftrees}{\mathcal{S}}
\newcommand{\later}{\mathord{\blacktriangleright}}
\newcommand{\Str}{\mathrm{Str}}
% IPsec notation
\newcommand{\Bytes}{\ensuremath{\mathrm{Bytes}}}
\newcommand{\SA}{\ensuremath{\mathrm{SA}}}
\newcommand{\SAD}{\ensuremath{\mathrm{SAD}}}
\newcommand{\SPD}{\ensuremath{\mathrm{SPD}}}
\newcommand{\SPI}{\ensuremath{\mathrm{SPI}}}
\newcommand{\ESP}{\ensuremath{\mathrm{ESP}}}
\newcommand{\AH}{\ensuremath{\mathrm{AH}}}
\newcommand{\IKE}{\ensuremath{\mathrm{IKEv2}}}
\newcommand{\Enc}{\ensuremath{\mathrm{Enc}}}
\newcommand{\Dec}{\ensuremath{\mathrm{Dec}}}
\newcommand{\MAC}{\ensuremath{\mathrm{MAC}}}

\newcommand{\redt}[1] {
	\textcolor{red}{#1}}

\newcommand{\oranget}[1] {
	\textcolor{orange}{#1}}

\newcommand{\purplet}[1] {
	\textcolor{purple}{#1}}

%%%%%%%%%%%%%%%%%%%%%%% Massangaben fuer die Arbeit %%%%%%%%%%%%%%%
\setlength{\textwidth}{15cm}

\setlength{\oddsidemargin}{35mm}
\setlength{\evensidemargin}{25mm}

\addtolength{\oddsidemargin}{-1in}
\addtolength{\evensidemargin}{-1in}

\ifthenelse{\equal{\biber}{true}}{\addbibresource{references.bib}}{}

%\makeindex
